% !TeX root = ../main.tex

\chapter{Divisibility}

\section{Prime Factorisation as Operator Factorisation}
% =================================================

Suppose

\[
n
=
p_1^{\alpha_1}
p_2^{\alpha_2}
\cdots
p_r^{\alpha_r}
\]

is the prime factorisation of \(n\).

Because

\[
T_mT_n=T_{mn},
\]

we obtain

\[
\boxed{
T_n
=
T_{p_1}^{\alpha_1}
T_{p_2}^{\alpha_2}
\cdots
T_{p_r}^{\alpha_r}.
}
\]

For example,

\[
12=2^2\cdot3,
\]

and therefore

\[
\boxed{
T_{12}=T_2^2T_3.
}
\]

Acting on \(\phi_1\),

\[
T_{12}\phi_1=\phi_{12}.
\]

But we may reach the same block through

\[
T_2T_2T_3\phi_1.
\]

Thus prime factorisation of integers becomes factorisation of
composition operators.

This suggests that the operators

\[
T_2,T_3,T_5,T_7,\ldots
\]

play the role of fundamental generators.

% =================================================
\section{Divisibility as Reachability}
% =================================================

We now obtain a direct interpretation of divisibility.

Suppose \(d,n\in\N\).

If

\[
d\mid n,
\]

then there exists an integer \(m\) such that

\[
n=dm.
\]

Consequently,

\[
T_m\phi_d
=
\phi_{dm}
=
\phi_n.
\]

Therefore

\[
\boxed{
d\mid n
\quad\Longleftrightarrow\quad
\phi_n
\text{ can be reached from }
\phi_d
\text{ by some }T_m.
}
\]

This allows us to interpret the divisibility relation geometrically as
a network of function transformations.

For example,

\[
\phi_2
\longrightarrow
\phi_4
\longrightarrow
\phi_8
\]

corresponds to

\[
2\mid4\mid8.
\]

Similarly,

\[
\phi_3
\longrightarrow
\phi_6
\longrightarrow
\phi_{12}
\]

corresponds to

\[
3\mid6\mid12.
\]

\subsection*{Why divisibility as a network of function transformations matters}

Interpreting the divisibility relation \(d\mid n\) as a directed edge
\(\phi_d\to\phi_n\) in a graph of functions turns an arithmetic relation
into a geometric and dynamical object. This viewpoint is useful for several reasons.

\paragraph{Visualization and intuition.}
Representing indices as nodes and divisibility as arrows makes chains and
reachability immediate. For instance,


\[
\phi_2\longrightarrow\phi_4\longrightarrow\phi_8
\]


visually encodes \(2\mid4\mid8\).

\paragraph{Operator algebra.}
The substitution operators \(T_m\) defined by \(T_mF(x)=F(x^m)\) satisfy
\(T_mT_n=T_{mn}\). Thus multiplicative identities among integers become
algebraic identities among operators, allowing analytic techniques to
study arithmetic phenomena.

\paragraph{Bridge between analysis and number theory.}
The same family \(\{\phi_n\}\) links functional analysis (composition,
differentiation, operator theory) with multiplicative number theory
(Dirichlet convolution, Möbius inversion). For example, operators of the
form \(\sum_{m\ge1}d_mT_m\) act on coefficient sequences by Dirichlet
convolution.

\paragraph{New invariants and classifications.}
Graph-theoretic notions such as distance, minimal predecessors, and
connected components yield invariants of indices that are not obvious in
the purely arithmetic setting; these invariants can detect factorization
patterns and structural symmetries.

\paragraph{Constructive and inverse methods.}
Natural operators (e.g. the divisor-sum operator \(D=\sum T_m\)) and
their inverses (via Möbius inversion) provide operational tools to
decompose and reconstruct series expressed in the \(\phi_n\) basis.

\paragraph{Potential applications.}
This perspective suggests new approaches in asymptotic analysis,
differential equations, combinatorics, and probabilistic models where
reachability by multiplication plays a role.


% =================================================
\section{A Divisibility Diagram}
% =================================================

The relation can be visualised using a small portion of the
divisibility lattice.

\begin{center}

\begin{tikzpicture}[
    node distance=1.5cm and 1.7cm,
    every node/.style={draw,rounded corners,inner sep=5pt},
    >=Stealth
]

\node (p1) {$\phi_1$};

\node (p2) [below left=of p1] {$\phi_2$};
\node (p3) [below right=of p1] {$\phi_3$};

\node (p4) [below=of p2] {$\phi_4$};
\node (p6) [below right=of p2] {$\phi_6$};

\node (p12) [below right=of p4] {$\phi_{12}$};

\draw[->] (p1) -- node[left,draw=none] {$T_2$} (p2);
\draw[->] (p1) -- node[right,draw=none] {$T_3$} (p3);

\draw[->] (p2) -- node[left,draw=none] {$T_2$} (p4);
\draw[->] (p2) -- node[above,draw=none] {$T_3$} (p6);

\draw[->] (p3) -- node[right,draw=none] {$T_2$} (p6);

\draw[->] (p4) -- node[left,draw=none] {$T_3$} (p12);
\draw[->] (p6) -- node[right,draw=none] {$T_2$} (p12);

\end{tikzpicture}

\end{center}

Different paths may lead to the same function because multiplication
is commutative.

For example,

\[
T_2T_3\phi_1
=
T_3T_2\phi_1
=
\phi_6.
\]

% =================================================
\section{Greatest Common Divisors and Least Common Multiples}
% =================================================

The divisibility interpretation also introduces naturally the
greatest common divisor and least common multiple.

Take two blocks

\[
\phi_m
\qquad\text{and}\qquad
\phi_n.
\]

A block \(\phi_d\) is a common predecessor when

\[
d\mid m
\qquad\text{and}\qquad
d\mid n.
\]

The largest possible such index is

\[
d=\gcd(m,n).
\]

Thus

\[
\boxed{
\phi_{\gcd(m,n)}
}
\]

is the greatest common predecessor of \(\phi_m\) and \(\phi_n\).

Similarly, a block \(\phi_k\) is a common descendant if

\[
m\mid k
\qquad\text{and}\qquad
n\mid k.
\]

The smallest such index is

\[
k=\operatorname{lcm}(m,n).
\]

Hence

\[
\boxed{
\phi_{\operatorname{lcm}(m,n)}
}
\]

is the smallest common descendant.

For example,

\[
\gcd(4,6)=2
\]

and

\[
\operatorname{lcm}(4,6)=12.
\]

Therefore

\[
\phi_2
\]

is the greatest common predecessor of

\[
\phi_4,\phi_6,
\]

while

\[
\phi_{12}
\]

is their smallest common descendant.

% =================================================
\section{Prime Blocks are Indecomposable}
% =================================================

Consider

\[
\phi_6(x)=b^{x^6}.
\]

Because

\[
6=2\cdot3,
\]

we can write

\[
\phi_6(x)
=
\phi_2(x^3)
=
\phi_3(x^2).
\]

Similarly,

\[
12=2\cdot6=3\cdot4,
\]

so

\[
\phi_{12}(x)
=
\phi_2(x^6)
=
\phi_3(x^4)
=
\phi_4(x^3)
=
\phi_6(x^2).
\]

However, suppose \(p\) is prime.

If

\[
\phi_p(x)=\phi_m(x^n)
\]

with

\[
m>1,
\qquad
n>1,
\]

then we would have

\[
p=mn,
\]

contradicting the primality of \(p\).

Therefore prime-indexed blocks are indecomposable under non-trivial
power composition.

\begin{proposition}

If \(p\) is prime, then there do not exist \(m,n>1\) such that

\[
\phi_p(x)=\phi_m(x^n).
\]

\end{proposition}

% =================================================
